\section{Results}\label{sec:results}

We estimate equation~\eqref{eq:stacked} separately for each of the five principal outcomes. Table~\ref{tab:main} reports the post-revision contrasts $\Delta_m$, and Figure~\ref{fig:event} shows the event-time path for licensing, the outcome with the largest and most precisely estimated contrast.

For licensing, the three pre-period coefficients are positive but imprecise: 0.30, 0.14, and 0.13 log points in event years $-4$ through $-2$. The joint pre-trend test does not reject that they are zero ($p=0.32$). The differential is 0.31 in the revision year, stays near 0.31 in years $+1$ and $+2$, and grows in years $+3$ through $+5$. Averaged over event years 0--5, the upward-minus-downward licensing gap is 0.495 log points (SE 0.140). Across 100,000 direction-label permutations, the two-sided nominal $p$-value is 0.032.

\begin{figure}[htbp]\centering
\includegraphics[width=0.82\textwidth]{figures/fig1_event_study_expanded.pdf}
\caption{Licensing around documentary-reviewed IP-policy revisions}\label{fig:event}
\begin{minipage}{0.9\textwidth}\footnotesize\textit{Notes:} Coefficients are the upward-minus-downward event-time differences $2\theta^{\text{lic}}_k$ from equation~\eqref{eq:stacked} for the preferred 34-event sample, with institution-clustered 95\% confidence intervals. Year $-1$ is the reference period.
\end{minipage}
\end{figure}

\begin{table}[htbp]\centering
\caption{Technology-transfer outcomes around IP-policy revisions}\label{tab:main}
\begin{threeparttable}\footnotesize
\resizebox{0.95\textwidth}{!}{%
\begin{tabular}{lrrrrrrrr}\toprule
 & \multicolumn{5}{c}{Upward minus downward} & \multicolumn{2}{c}{Relative to controls} & \\\cmidrule(lr){2-6}\cmidrule(lr){7-8}
Outcome & Gap & SE & Perm.\ $p$ & BH $q$ & Pre $p$ & Upward & Downward & $N$ \\\midrule
Log licenses and options & 0.495 & (0.140) & 0.032 & 0.159 & 0.32 & 0.279 & -0.216 & 26,555 \\
Filing margin & 0.392 & (0.155) & 0.097 & 0.243 & 0.75 & 0.143 & -0.249 & 26,679 \\
Log new patent applications & 0.260 & (0.175) & 0.205 & 0.341 & 0.89 & 0.136 & -0.124 & 26,679 \\
Log U.S. patents issued & 0.137 & (0.190) & 0.541 & 0.541 & 0.47 & 0.171 & 0.034 & 26,151 \\
Log invention disclosures & -0.132 & (0.095) & 0.380 & 0.475 & 0.40 & -0.007 & 0.124 & 26,679 \\
\bottomrule
\end{tabular}}
\begin{tablenotes}[flushleft]\footnotesize
\item \textit{Notes:} Preferred sample: 34 documentary-reviewed revisions, eight upward and 26 downward; target window $-4$ to $+5$, with year $-1$ omitted. Fixed effects: stack $\times$ institution and stack $\times$ calendar year $\times$ research-size-tercile $\times$ public/private; control: log research expenditure. Permutation $p$: two-sided 100,000-draw reassignment test preserving eight upward labels. BH $q$: Benjamini--Hochberg adjustment across the five outcomes. Pre $p$: joint Wald test of the three pre-period direction coefficients. $N$ counts stacked institution-year-event observations after outcome-specific missing-value restrictions; it is not the number of unique institution-years or policy events. Permutation $p$- and $q$-values are Monte Carlo estimates (simulation standard error about 0.0006 for the licensing $p$-value), so their third decimal can differ slightly across computing platforms.
\end{tablenotes}
\end{threeparttable}
\end{table}

The licensing differential comes from movement in both directions. Relative to non-revising controls, the post-period average is $+0.279$ after upward revisions and $-0.216$ after downward ones. The two are not the same size, but neither direction produces the gap on its own. Under the stricter post-support rule, the licensing gap is 0.499 (SE 0.133), with $+0.230$ for upward and $-0.269$ for downward revisions.

The other outcomes help locate the pattern; they do not extend it. The filing margin has a positive gap of 0.392 and no detectable pre-period difference, but its nominal permutation $p$-value is 0.097. New patent applications have a positive gap of 0.260 ($p=0.205$), and patents issued show a smaller positive differential. For invention disclosures, the gap is negative and imprecise. We see none of the increase that would be expected if the main post-revision pattern ran through more faculty entering the disclosure process.

Accounting for multiple tests changes how strong this evidence looks. Licensing is unusual on its own in the permutation distribution, but the Benjamini--Hochberg adjustment across the five outcomes gives $q\approx0.16$, which does not meet conventional false-discovery-rate thresholds. We read licensing as the strongest outcome-specific pattern in the data, not as evidence that revisions shift the whole technology-transfer pipeline.

Where the pattern appears is still informative. Licensing depends on what happens after disclosure: evaluation, the search for a licensee, negotiation, and continued cooperation among inventors, TTO staff, and firms. The filing margin and new patent applications also sit downstream of disclosure and move the same way, though too imprecisely to support separate conclusions. With no matching rise in disclosures, the simple story that more faculty come to the TTO after a revision fits the data poorly.

\subsection*{Robustness to the preferred sample definition}

The move from the earlier 25-event set to the preferred sample came from documentary review, not from looking at outcome estimates. We reviewed the 37 mechanically eligible pairs that had not been coded by hand under the same comparability and timing rules, and nine more events passed. Table~\ref{tab:robust} compares the preferred sample with the stricter post-support sample and with the original 25-event set.

\begin{table}[htbp]\centering
\caption{Alternative documentary event samples: licensing and filing margin}\label{tab:robust}
\footnotesize
\resizebox{0.96\textwidth}{!}{%
\begin{tabular}{lcrrrrrrr}\toprule
 & & \multicolumn{4}{c}{Log licenses and options} & \multicolumn{3}{c}{Filing margin} \\\cmidrule(lr){3-6}\cmidrule(lr){7-9}
Sample & Up/down & Gap & SE & Perm.\ $p$ & Pre $p$ & Gap & SE & Perm.\ $p$ \\\midrule
Preferred documentary sample & 8/26 & 0.495 & (0.140) & 0.032 & 0.32 & 0.392 & (0.155) & 0.097 \\
Strict post-support sample & 7/22 & 0.499 & (0.133) & 0.026 & 0.52 & 0.361 & (0.158) & 0.119 \\
Earlier 25-event set & 6/19 & 0.621 & (0.136) & 0.009 & 0.62 & 0.380 & (0.170) & 0.153 \\
\bottomrule
\end{tabular}}
\vspace{2pt}

\parbox{0.94\textwidth}{\footnotesize\textit{Notes:} The preferred and strict samples use 100,000 Monte Carlo direction assignments. The earlier 25-event row uses exact enumeration over all $\binom{25}{6}=177{,}100$ direction assignments.}
\end{table}

The stricter sample drops preferred-sample events that have no observed year after the revision, and it applies that rule to the original and the newly added events alike. Its licensing estimate is nearly the same as the preferred one, with a weaker pre-period differential. The result does not come from treating the last AUTM year as if it offered a long post-revision window.

No single revision drives the licensing estimate. Dropping each of the 34 events in turn gives gaps between 0.436 and 0.559 (Appendix Figure~\ref{fig:loo}), and the estimate stays positive whichever of the eight upward or 26 downward events is left out.

\subsection*{Pre-revision balance and selection}

Appendix Table~\ref{tab:balance} compares upward and downward events before the revision. Their pre-revision licensing trends are essentially the same. Differences in pre-revision licensing, research expenditure, public or private control, and documentary classification are not unusual under label permutations. Upward revisions come somewhat later and at institutions with more disclosures, but these differences are imprecise.

The clearest imbalance is in the predecessor policy. Mean PCSI of the earlier document is 0.384 for upward revisions and 0.485 for downward revisions, a standardized difference of $-1.65$ with permutation $p=0.002$. This is a real concern about selection and mean reversion, not a small balance problem. Universities whose earlier policies were relatively restrictive have more room to move upward, and their reasons for doing so may be tied to wider organizational change. This imbalance is a main reason we do not give the direction contrast a causal reading, even though the licensing pre-trends look alike.

\subsection*{Revision content and supporting annual evidence}

Across all 126 threshold revisions, downward revisions are more likely than upward ones to add conflict-of-interest language (Table~\ref{tab:content}). The comparison is based on keyword presence, not on a legal reading of what the provisions require. Provision-level hand coding exists only for the original 25 content-coded events (Appendix~\ref{app:provisions}). We report it as supplementary evidence and did not extend those labels to the newly added events after seeing the expanded-sample results.

\begin{table}[htbp]\centering
\caption{Content added in upward and downward revisions}\label{tab:content}
\begin{threeparttable}\small\begin{tabular}{lrrr}\toprule
 & Upward & Downward & Fisher $p$ \\\midrule
\multicolumn{4}{l}{\textit{Original 25-event content-coded subset} (6 upward, 19 downward)} \\
\quad Adds present assignment language & 33\% & 16\% & 0.562 \\
\quad Adds equity language & 33\% & 47\% & 0.661 \\
\quad Adds conflict of interest language & 0\% & 37\% & 0.137 \\
\quad Adds startup language & 17\% & 11\% & 1.000 \\
\quad Adds copyright language & 33\% & 32\% & 1.000 \\
\multicolumn{4}{l}{\textit{All revisions} (46 upward, 80 downward)} \\
\quad Adds present assignment language & 9\% & 16\% & 0.287 \\
\quad Adds equity language & 28\% & 35\% & 0.554 \\
\quad Adds conflict of interest language & 13\% & 32\% & 0.019 \\
\quad Adds startup language & 15\% & 16\% & 1.000 \\
\quad Adds copyright language & 13\% & 20\% & 0.465 \\
\bottomrule\end{tabular}
\begin{tablenotes}[flushleft]\footnotesize
\item \textit{Notes:} Share of revisions whose later document contains language absent from the earlier document, detected by keyword search: present assignment (``hereby assign''), equity or stock, conflict(s) of interest, startup/spin-off/new venture, and copyright. Two-sided Fisher exact $p$-values are shown. Keyword presence is descriptive and does not establish the legal content or administrative effect of a provision.
\end{tablenotes}\end{threeparttable}\end{table}

The annual panel points the same general way, with much less precision. Table~\ref{tab:twfe} reports regressions of each outcome on contemporaneous and lagged PCSI. The filing-margin coefficient is positive at every horizon, the licensing coefficient is positive for several lags but imprecise, and the disclosure coefficients are negative at every lag shown. No lag-one coefficient survives the adjustment across the six annual-panel outcomes. These models look at the data differently, but they cannot replace the revision design: most annual PCSI values are carried forward, and the annual panel does not apply the documentary timing and comparability rules.

\begin{table}[htbp]\centering
\caption{Supporting evidence: continuous PCSI in two-way fixed-effects models}\label{tab:twfe}
\begin{threeparttable}\footnotesize\setlength{\tabcolsep}{3pt}\begin{tabular}{lrrrrrrr}\toprule
Outcome & Contemp. & Lag 1 & Lag 2 & Lag 3 & Lag 4 & Lag 5 & BH $q$, lag 1 \\\midrule
Invention disclosures & -0.319 & -0.368 & -0.306 & -0.343 & -0.316 & -0.102 & 0.345 \\
 & (0.375) & (0.374) & (0.367) & (0.374) & (0.391) & (0.379) & \\
Filing margin & 0.647 & 0.938 & 1.200 & 0.974 & 0.510 & 0.458 & 0.345 \\
 & (0.527) & (0.505) & (0.520) & (0.468) & (0.445) & (0.496) & \\
New patent applications & 0.328 & 0.569 & 0.894 & 0.631 & 0.195 & 0.356 & 0.345 \\
 & (0.560) & (0.523) & (0.537) & (0.535) & (0.477) & (0.484) & \\
Total patent applications & 0.516 & 0.589 & 0.640 & 0.316 & 0.259 & 0.237 & 0.345 \\
 & (0.573) & (0.525) & (0.499) & (0.487) & (0.498) & (0.468) & \\
Patents issued & 0.266 & 0.548 & 1.205 & 0.981 & 0.774 & 1.300 & 0.345 \\
 & (0.560) & (0.577) & (0.634) & (0.573) & (0.631) & (0.578) & \\
Licenses/options & 0.518 & 0.596 & 0.400 & 0.274 & 0.161 & -0.451 & 0.345 \\
 & (0.490) & (0.452) & (0.406) & (0.441) & (0.467) & (0.503) & \\
\bottomrule\end{tabular}
\begin{tablenotes}[flushleft]\footnotesize
\item \textit{Notes:} Coefficients on PCSI (contemporaneous or lagged $k$ calendar years); institution-clustered standard errors in parentheses. Models include institution fixed effects for the 134 harmonized institutions with observed PCSI, year fixed effects, and lagged log research expenditure, log licensing FTEs, and inventor royalty share. Filing-margin samples are $N=2,339, 2,271, 2,132, 2,005, 1,891, 1,772$ from contemporaneous to lag 5. BH $q$: Benjamini--Hochberg adjustment across the six lag-one outcomes shown.
\end{tablenotes}\end{threeparttable}\end{table}
